\section*{Capítulo \ref{ch:b1:weathering-clays} --- Intemperismo e argilas}

\begin{solution}{exo:b1:weathering-clays:1}
\ce{K+}: $1/1.38 = 0.72$; \ce{Fe^2+}: $2/0.78 = 2.56$; ambos abaixo de 3, cátions solúveis. \ce{Fe^3+}: $3/0.645 = 4.65$; \ce{Ti^4+}: $4/0.605 = 6.61$; ambos entre 3 e 10, hidrolisados. Em um solo encharcado, sem oxigênio, o ferro é liberado como \ce{Fe^2+}, que permanece em solução e é levado pela água; em um solo bem drenado, o oxigênio o transforma em \ce{Fe^3+}, que precipita de imediato como óxidos e hidróxidos e permanece.
\end{solution}

\begin{solution}{exo:b1:weathering-clays:2}
A incha com glicol e colapsa ao aquecer: uma \omterm{def:b1:weathering-clays:clay-mineral}{esmectita}. B está em
\SI{7}{\angstrom} e é destruída a \SI{550}{\celsius}: \omterm{def:b1:weathering-clays:clay-mineral}{caulinita}. C, em \SI{10}{\angstrom} qualquer que seja o tratamento: \omterm{def:b1:weathering-clays:clay-mineral}{ilita}. D, em \SI{14}{\angstrom}, sem inchamento nem colapso: \omterm{def:b1:weathering-clays:clay-mineral}{clorita}.
\end{solution}

\begin{solution}{exo:b1:weathering-clays:3}
Dois ortoclásios trazem os 2 alumínios de uma \omterm{def:b1:weathering-clays:clay-mineral}{caulinita}; $6 - 2 = 4$ silícios saem como \ce{H4SiO4}; 2 \ce{K+} saem, balanceados por 2 \ce{H+}; oxigênio, $16 + w = 9 + 16$, dá $w = 9$:
\[
  \ce{2 KAlSi3O8 + 2 H+ + 9 H2O -> Al2Si2O5(OH)4 + 2 K+ + 4 H4SiO4}.
\]
O hidrogênio confere: $2 + 18 = 20 = 4 + 16$. Todo o potássio sai, e $4/6 =
2/3$ do silício.
\end{solution}

\begin{solution}{exo:b1:weathering-clays:4}
(a) A massa não muda: $10 \times 999.87/916.72 = \qty{10.91}{cm^3}$.
(b) Entre cerca de \qty{-4}{\celsius} e \qty{-15}{\celsius}. Logo abaixo de zero o gelo não consegue acumular pressão suficiente para cunhar a fissura; em uma geada mais forte a água não se move mais pelos poros em direção ao gelo, que deixa de crescer. (c) Pela
\cref{prop:b1:weathering-clays:surface} com $n = \qty{1}{m}/\qty{1}{cm} = 100$,
a área é multiplicada por 100: de \qty{6}{m^2} a \qty{600}{m^2}.
\end{solution}

\begin{solution}{exo:b1:weathering-clays:5}
$S/m = 6/(\rho d)$: para \qty{1}{cm}, $6/(2600\times0.01) =
\qty{0.231}{m^2.kg^{-1}}$; para \qty{1}{mm}, \qty{2.31}{m^2.kg^{-1}}; para
\qty{1}{\micro m}, \qty{2308}{m^2.kg^{-1}}, ou \qty{2.3}{m^2.g^{-1}}. As reações químicas ocorrem na superfície dos grãos, de modo que sua taxa por quilograma de rocha cresce com a superfície específica, dez mil vezes do centímetro ao micrômetro. O \omterm{def:b1:weathering-clays:weathering}{intemperismo físico} corta a rocha em pedaços cada vez menores e abre fissuras pelas quais a água entra: ele multiplica a superfície sobre a qual o \omterm{def:b1:weathering-clays:weathering}{intemperismo químico} pode agir.
\end{solution}

\begin{solution}{exo:b1:weathering-clays:6}
Uma tonelada de albita é $\num{e6}/262.22 = \qty{3813.6}{mol}$. Pela
\cref{prop:b1:weathering-clays:feldspars}, dois moles dão um de \omterm{def:b1:weathering-clays:clay-mineral}{caulinita}:
$\qty{1906.8}{mol} \times \qty{258.16}{g.mol^{-1}} = \qty{492.3}{kg}$. Dois dos três átomos de silício saem: $2\times3813.6 = \qty{7627.2}{mol}$ de
\ce{SiO2}, \qty{458.2}{kg}. Todo o sódio sai: $3813.6\times22.99 =
\qty{87.7}{kg}$. Menos da metade da massa do \omterm{def:b1:rock-forming-minerals:groups}{feldspato} permanece como argila.
\end{solution}

\begin{solution}{exo:b1:weathering-clays:7}
Pela \cref{prop:b1:weathering-clays:immobile}, um quilograma de granito deixa
$0.30/0.60 = \qty{0.5}{kg}$ de saprólito. Fração perdida: sódio $1 -
(0.35/3.5)(0.5) = 0.95$; silício $1 - (60/72)(0.5) = 0.583$; alumínio $1 -
(27/14)(0.5) = 0.036$. O alumínio é praticamente retido, o sódio quase se foi por completo, e quase \qty{60}{\percent} do silício também, embora grande parte do que resta esteja retida no quartzo: os \omterm{def:b1:rock-forming-minerals:groups}{feldspatos} estão se transformando em \omterm{def:b1:weathering-clays:clay-mineral}{caulinita}, \omterm{def:b1:weathering-clays:hydrolysis-degree}{monossialitização}.
\end{solution}

\begin{solution}{exo:b1:weathering-clays:8}
Pirofilita: $x = y = 0$, carga 0. Camada de moscovita: $x = 1$, carga $-1$.
Camada de beidelita: $x = 0.5$, carga $-0.5$. Camada de montmorilonita: $y = 0.33$,
carga $-0.33$. As duas camadas de \omterm{def:b1:weathering-clays:clay-mineral}{esmectita}, beidelita e montmorilonita, podem inchar: sua pequena carga é balanceada por poucos cátions entre as camadas, que as mantêm frouxamente e atraem água à medida que se hidratam. A pirofilita não tem cátions entre as camadas para se hidratar, e a carga $-1$ da moscovita é balanceada por íons potássio que travam as camadas.
\end{solution}

\begin{solution}{exo:b1:weathering-clays:9}
(a) Mg: 2 e 2; Si: 1 e 1; C: 4 e 4; O: $4 + 8 + 4 = 16$ e $12 + 4 =
16$; H: 8 e $4 + 4 = 8$; carga: 0 e $4 - 4 = 0$. (b) A massa molar da forsterita é $2\times24.305 + 28.085 + 4\times15.999 =
\qty{140.69}{g.mol^{-1}}$, logo uma tonelada são \qty{7108}{mol}. Ela consome $4\times7108 = \qty{28431}{mol}$ de \ce{CO2}, \qty{1251}{kg}; com dois magnésios por unidade de fórmula, \cref{prop:b1:weathering-clays:co2} armazena $2\times7108$ moles, \qty{626}{kg}, metade da quantidade consumida. (c) Uma tonelada de calcita, \qty{9991}{mol}, consome \qty{9991}{mol} de \ce{CO2} (\qty{440}{kg}) ao se dissolver, e devolve todo ele quando o carbonato precipita de novo: nada é armazenado.
\end{solution}

\begin{solution}{exo:b1:weathering-clays:10}
(a) $3 < 3/r < 10$ dá $\qty{0.3}{\angstrom} < r < \qty{1.0}{\angstrom}$. (b) $3/0.27 = 11.1 > 10$: o boro viaja como oxiânion, borato. (c) $2/0.83 =
2.41$, um cátion solúvel; $4/0.53 = 7.55$, um hidrolisado. Em água pobre em oxigênio o manganês se dissolve como \ce{Mn^2+}; onde a nascente encontra o ar, ele é oxidado a \ce{Mn^4+}, que precipita de imediato como um óxido preto sobre a rocha. (d) Seu \omterm{def:b1:weathering-clays:ionic-potential}{potencial iônico}, $4/0.40 = 10.0$, jaz na fronteira entre hidrolisados e oxiânions: o silício não é retido como o alumínio nem levado embora como o enxofre, e a fração lixiviada depende de quanta água passa, base dos três graus de \omterm{def:b1:weathering-clays:hydrolysis}{hidrólise}.
\end{solution}

\begin{solution}{exo:b1:weathering-clays:11}
(a) Com \qty{12.011}{g.mol^{-1}} de carbono, \num{11.7e12} a
\qty{13.6e12}{mol} é 0.141 a 0.163 bilhão de toneladas de carbono por ano; com \qty{44.009}{g.mol^{-1}}, 0.515 a 0.599 bilhão de toneladas de gás carbônico. (b) $4.08/13.6 = \qty{30}{\percent}$ a $4.08/11.7 = \qty{35}{\percent}$.
(c) Metade do gás carbônico consumido é armazenada quando o cálcio e o magnésio precipitam como carbonato: 0.070 a 0.082 bilhão de toneladas de carbono por ano. O valor real é menor, porque parte do bicarbonato sai com sódio e potássio, que não formam carbonato no fundo do mar.
\end{solution}

\begin{solution}{exo:b1:weathering-clays:12}
Si/Al $= (w_{\ce{SiO2}}/60.08)/(2w_{\ce{Al2O3}}/101.96)$. A: $0.7748/0.7748 =
1.00$, a razão da \omterm{def:b1:weathering-clays:clay-mineral}{caulinita}, \omterm{def:b1:weathering-clays:hydrolysis-degree}{monossialitização}. B: $0.9987/0.4904 = 2.04$, um \omterm{def:b1:weathering-clays:clay-mineral}{argilomineral} 2:1, \omterm{def:b1:weathering-clays:hydrolysis-degree}{bissialitização}. C: $0.0333/1.2358 = 0.027$, quase nenhum silício restante, gibbsita, \omterm{def:b1:weathering-clays:hydrolysis-degree}{alitização}. Da lixiviação menos à mais intensa, isto é, em direção ao clima mais quente e mais úmido: B, A, C.
\end{solution}

\begin{solution}{pb:b1:weathering-clays:1}
\textbf{1.} \ce{2 NaAlSi3O8 + 2 H+ + 9 H2O -> Al2Si2O5(OH)4 + 2 Na+ + 4 H4SiO4}.

\textbf{2.} \ce{CaAl2Si2O8 + 2 H+ + H2O -> Al2Si2O5(OH)4 + Ca^2+}. Nenhum silício
sai: a anortita já tem o silício por alumínio da \omterm{def:b1:weathering-clays:clay-mineral}{caulinita}.

\textbf{3.} Uma gibbsita toma um alumínio, metade de uma anortita: dobrando,
\ce{CaAl2Si2O8 + 2 H+ + 6 H2O -> 2 Al(OH)3 + Ca^2+ + 2 H4SiO4}. Oxigênio:
$8 + 6 = 6 + 8$; hidrogênio: $2 + 12 = 6 + 8$.

\textbf{4.} \ce{5 NaAlSi3O8 + 4 H+ + 16 H2O -> 2 Na_{0.5}Al2(Si_{3.5}Al_{0.5})O10(OH)2 + 4 Na+ + 8 H4SiO4}; % ce-unbalanced-ok: decimal subscripts the checker cannot read; balance tested in tests/bachelor-1/test_weathering_clays.py
hidrogênio, $4 + 32 = 36$ à esquerda e $2\times2 + 8\times4 = 36$ à
direita, confere.

\textbf{5.} Cada \ce{H+} vem com um \ce{HCO3-} de um \ce{CO2} e uma
\ce{H2O}:
\ce{2 NaAlSi3O8 + 2 CO2 + 11 H2O -> Al2Si2O5(OH)4 + 2 Na+ + 2 HCO3- + 4 H4SiO4}
e \ce{CaAl2Si2O8 + 2 CO2 + 3 H2O -> Al2Si2O5(OH)4 + Ca^2+ + 2 HCO3-}.

\textbf{6.} A albita para beidelita é \omterm{def:b1:weathering-clays:hydrolysis-degree}{bissialitização}; albita e anortita para
\omterm{def:b1:weathering-clays:clay-mineral}{caulinita} são \omterm{def:b1:weathering-clays:hydrolysis-degree}{monossialitização}; anortita para gibbsita é \omterm{def:b1:weathering-clays:hydrolysis-degree}{alitização}.

\textbf{7.} $\num{e6}/262.22 = \qty{3813.6}{mol}$ de albita;
$\num{e6}/278.20 = \qty{3594.5}{mol}$ de anortita.

\textbf{8.} Albita: $3813.6/2 = \qty{1906.8}{mol}$ de \omterm{def:b1:weathering-clays:clay-mineral}{caulinita},
\qty{492.3}{kg}. Anortita: \qty{3594.5}{mol}, \qty{928.0}{kg}.

\textbf{9.} Albita: $2\times3813.6 = \qty{7627.2}{mol}$ de \ce{SiO2},
\qty{458.2}{kg}, e \qty{3813.6}{mol} de sódio, \qty{87.7}{kg}. Anortita:
sem sílica, e \qty{3594.5}{mol} de cálcio, \qty{144.1}{kg}.

\textbf{10.} Albita: $1/2.62 = \qty{0.382}{m^3}$; sua \omterm{def:b1:weathering-clays:clay-mineral}{caulinita},
$0.4923/2.6 = \qty{0.189}{m^3}$: o sólido perde \qty{50.4}{\percent} de
seu volume. Anortita: $1/2.75 = \qty{0.364}{m^3}$; sua \omterm{def:b1:weathering-clays:clay-mineral}{caulinita},
$0.9280/2.6 = \qty{0.357}{m^3}$: uma perda de apenas \qty{1.8}{\percent}.

\textbf{11.} O espaço deixado torna-se poros, preenchidos por água e ar: um
saprólito sobre um granito rico em albita é leve, poroso e friável, seus grãos de quartzo soltos em um pouco de argila. A anortita se transforma em quase seu próprio volume de \omterm{def:b1:weathering-clays:clay-mineral}{caulinita}, de modo que o gabro intemperizado mantém a maior parte de sua matéria sólida.

\textbf{12.} Sódio ($z/r = 0.98$) e cálcio (2.0) são cátions solúveis e saem; alumínio (5.6) é um hidrolisado e permanece todo na \omterm{def:b1:weathering-clays:clay-mineral}{caulinita}; silício (10.0), na fronteira, sai em parte da albita e não sai da anortita, cujo silício a \omterm{def:b1:weathering-clays:clay-mineral}{caulinita} pode reter.

\textbf{13.} Um \ce{CO2} por albita: \qty{3813.6}{mol}, $3813.6 \times 44.01 =
\qty{167.8}{kg}$.

\textbf{14.} Dois por anortita: \qty{7189.0}{mol}, \qty{316.4}{kg}.

\textbf{15.} $7189.0 \times 61.02 = \qty{438.7}{kg}$ de bicarbonato.

\textbf{16.} Os organismos o precipitam como carbonato de cálcio,
\ce{Ca^2+ + 2 HCO3- -> CaCO3 + CO2 + H2O}: para cada cálcio, um dos dois
\ce{CO2} volta ao ar, \qty{3594.5}{mol} ou \qty{158.2}{kg}, e o outro fica enterrado no carbonato.

\textbf{17.} O sódio não forma carbonato no fundo do mar: seu bicarbonato não pode ser armazenado como rocha carbonática, e o armazenamento de longo prazo é contado só a partir de cálcio e magnésio.

\textbf{18.} Uma tonelada de calcita consome \qty{9991}{mol} de \ce{CO2} quando se dissolve e devolve tudo quando precipita: nada é armazenado, e somente o intemperismo de silicatos de cálcio e magnésio retira gás carbônico do ar para sempre.

\textbf{19.} $\qty{200}{m^3.s^{-1}} \times \qty{0.3}{mol.m^{-3}} =
\qty{60}{mol.s^{-1}}$; em um ano de \qty{3.156e7}{s}, \qty{1.893e9}{mol}, ou
$\num{1.893e9}/\num{50000} = \qty{37870}{mol}$ por quilômetro quadrado.

\textbf{20.} Dois bicarbonatos por anortita: \qty{9.467e8}{mol} por ano,
$\num{9.467e8}\times\qty{278.20}{g} = \qty{263400}{t}$ por ano, ou
\qty{5.27}{t} por quilômetro quadrado.

\textbf{21.} A rocha fresca consumida é $5.27/0.20 = \qty{26.3}{t}$, ou
$\num{26340}/2700 = \qty{9.75}{m^3}$, por quilômetro quadrado e ano: uma camada de
\qty{9.75}{\micro m} de espessura por ano, \qty{9.8}{m} por milhão de anos.

\textbf{22.} Um \ce{CO2} armazenado por cálcio: $\num{9.467e8} \times 44.01 =
\qty{4.17e10}{g}$, cerca de \qty{41700}{t} por ano.

\textbf{23.} Ele conta apenas a matéria levada embora em solução. A \omterm{def:b1:weathering-clays:clay-mineral}{caulinita} e o quartzo deixados pelo intemperismo também são removidos pela erosão, como lama e areia no rio, e não são contados; nem o intemperismo dos outros \omterm{def:b1:rock-forming-minerals:mineral}{minerais} da rocha. E se parte do bicarbonato vier de \omterm{def:b1:rock-forming-minerals:groups}{feldspatos} sódicos ou potássicos, que dão um bicarbonato por unidade de fórmula em vez de dois, mais rocha teve de se intemperizar para fornecê-lo.

\textbf{24.} Pela \cref{prop:b1:weathering-clays:co2}, um mol de \ce{CO2} é armazenado por mol de anortita: $3594.5 \times \qty{44.01}{g} = \qty{158.19}{kg}$. O intemperismo de uma tonelada de anortita retira para sempre
\textbf{158}~kg de gás carbônico.
\end{solution}
